\chapter{Cesarean Section}

\section*{Overview}
A cesarean section, or C-section, is delivery of a baby through surgical incisions in the abdomen and uterus.

\section*{Why This Test or Procedure Is Ordered}
It may be planned or performed urgently for placenta previa, fetal distress, abnormal presentation, labor complications, prior uterine surgery, or maternal or fetal conditions that make vaginal birth unsafe.

\section*{How This Test or Procedure Is Performed}
Regional anesthesia is usually used. The surgeon opens the abdomen and uterus, delivers the baby and placenta, and closes the uterine and abdominal incisions.

\section*{Risks and Recovery}
Infection, bleeding, blood clots, injury to nearby organs, anesthesia complications, and risks in future pregnancies are possible. Recovery is generally longer than after uncomplicated vaginal birth.

\section*{References}
\begin{enumerate}
  \item MedlinePlus. Cesarean Section. U.S. National Library of Medicine; 2025.
  \item Current Medical Diagnosis and Treatment. McGraw-Hill; 2025.
\end{enumerate}
\clearpage
